\chapter{Chronology}
\label{app:chronology}

A timeline of the results the book rests on, from Bezzel, Nauck and Gauss to the computations reported here.

\begin{description}[style=nextline]
\item[1848 to 1850] Bezzel poses the eight queens problem; Nauck publishes the 92 solutions and poses the $N$ queens problem; Gauss discusses the eight queens in letters to Schumacher.
\item[1874] Pauls gives explicit solutions for all $N \ge 4$.
\item[1900] Hilbert's Paris problems; the tenth asks for a decision procedure for Diophantine equations.
\item[1905] Lorenz introduces his curve for the concentration of wealth.
\item[1918] P\'olya shows toroidal solutions exist only when $\gcd(N, 6) = 1$.
\item[1928] Hilbert and Ackermann state the Entscheidungsproblem.
\item[1929 to 1934] Hardy, Littlewood and P\'olya characterise majorization.
\item[1931] G\"odel's incompleteness theorems.
\item[1936] Church and Turing show the Entscheidungsproblem is unsolvable; Turing's machine.
\item[1937] Shannon's master's thesis on switching circuits.
\item[1945] von Neumann's First Draft of a Report on the EDVAC.
\item[1946 to 1949] Ulam and von Neumann invent Monte Carlo; Metropolis names it; Metropolis and Ulam publish.
\item[1948 to 1950] Shannon's information theory, circuit counting bound, secrecy systems and chess paper.
\item[1951] von Neumann's rejection method and coin unbiasing; Kahn and Harris's splitting.
\item[1953] The Metropolis algorithm.
\item[1956] G\"odel's letter to von Neumann; von Neumann's probabilistic logics.
\item[1959] Rabin and Scott define nondeterminism.
\item[1965] Hartmanis and Stearns's time hierarchy; Cobham and Edmonds propose polynomial time.
\item[1966] Haber's stratified Monte Carlo quadrature.
\item[1970] Matiyasevich settles Hilbert's tenth problem.
\item[1971 to 1973] Cook, Karp and Levin establish $\NP$ completeness.
\item[1973] Peskun's ordering; Br\`egman's permanent inequality.
\item[1975] Baker, Gill and Solovay's relativization barrier; Ladner's theorem; Knuth's backtrack estimation.
\item[1976 to 1980] Solovay and Strassen, Miller and Rabin's primality tests; Gill's randomized classes.
\item[1978] Adleman's theorem $\RP \subseteq \Ppoly$ (extended to $\BPP$ by Bennett and Gill~\cite{BennettGill} in 1981); Yakowitz, Krimmel and Szidarovszky's weighted Monte Carlo integration.
\item[1979] Valiant defines $\sharpP$ and proves the permanent complete; Cook and Reckhow found proof complexity.
\item[1981] Diaconis and Shahshahani's random transpositions; van der Waerden's conjecture proved.
\item[1982 to 1984] Yao, Blum and Micali define pseudorandomness; Geman and Geman's Gibbs sampler; parity not in $\mathrm{AC}^0$.
\item[1983] Sipser, G\'acs and Lautemann place $\BPP$ in the hierarchy; Karp and Luby count DNF; Stockmeyer's approximate counting; Kirkpatrick, Gelatt and Vecchi's simulated annealing.
\item[1985] Razborov's monotone lower bound; Haken's resolution lower bound; interactive proofs.
\item[1986] Jerrum, Valiant and Vazirani relate counting and sampling; Valiant and Vazirani's isolation theorem; H{\aa}stad's switching lemma; Levin's average case completeness.
\item[1987 to 1988] Swendsen and Wang's cluster algorithm; Edwards and Sokal's representation.
\item[1989] Jerrum and Sinclair's FPRAS for the dense permanent and the Ising model.
\item[1990 to 1992] Gelfand and Smith's Gibbs sampling in statistics; Geyer's practical MCMC and the witch's hat discussion; Masry and Cambanis's trapezoidal Monte Carlo.
\item[1991] Dyer, Frieze and Kannan's volume algorithm; Toda's theorem; Berg and Neuhaus's multicanonical sampler; Applegate and Kannan on log concave sampling.
\item[1992] Shamir's $\mathrm{IP} = \mathrm{PSPACE}$; the PCP theorem.
\item[1994] Nisan and Wigderson's hardness versus randomness; Razborov and Rudich's natural proofs; Frieze, Kannan and Polson's log concave Metropolis walk; Tierney's convergence theory.
\item[1995] Diaconis and Holmes on the queens as Monte Carlo; Impagliazzo's five worlds; Kannan, Lov\'asz and Simonovits's conjecture; Hesselbo and Stinchcombe's $1/k$ ensemble.
\item[1996] Polson's convergence of MCMC; CDCL solvers.
\item[1997] Impagliazzo and Wigderson: $\Pc = \BPP$ under circuit lower bounds; Bubley and Dyer's path coupling.
\item[1998 to 1999] Andreev, Clementi and Rolim's hitting sets; Higdon's auxiliary variables; Gore and Jerrum's torpid Swendsen Wang; Damien, Wakefield and Walker's slice sampling.
\item[2000] Bellare, Goldreich and Petrank's uniform generation; Brassard, H{\o}yer, Mosca and Tapp's amplitude estimation; the Clay prize.
\item[2001] Wang and Landau's algorithm; Ben Sasson and Wigderson's width bound.
\item[2002] Agrawal, Kayal and Saxena: primality in $\Pc$; M\'ezard, Parisi and Zecchina's survey propagation.
\item[2004] Jerrum, Sinclair and Vigoda's FPRAS for the permanent; Kabanets and Impagliazzo's barrier; Szegedy's quantum walks.
\item[2006] Skilling's nested sampling; Lov\'asz and Vempala's simulated annealing for volume.
\item[2007 to 2008] Dinur's PCP proof; Achlioptas and Coja Oghlan's clustering transition; Aaronson and Wigderson's algebrization.
\item[2009] \v{S}tefankovi\v{c}, Vempala and Vigoda's adaptive annealing.
\item[2010 to 2011] Sly's computational transition; Williams's $\mathrm{ACC}^0$ lower bound; Chopin and Robert's nested sampling theory.
\item[2013] Birge, Chang and Polson's split sampling.
\item[2014] Gamarnik and Sudan's overlap gap property; Keevash's existence of designs; Polson and Scott's vertical likelihood Monte Carlo.
\item[2015 to 2016] Ding, Sly and Sun's satisfiability threshold; Babai's graph isomorphism; the exact count $Q(27)$ (Preusser and Engelhardt).
\item[2017] Gent, Jefferson and Nightingale: queens completion is $\NP$ complete.
\item[2021 to 2023] Simkin's theorem on $Q(N)$; Chen, Klartag and Lehec on KLS; Polson and Sokolov's generative Bayesian computation; Polson, Sokolov and Xu's quantum Bayesian computation.
\item[2024] Polson and Sokolov, Counting $N$ Queens.
\item[2026] This book: the main equivalence, the Yakowitz improvements, the giant component to $N = 17$, the mixing exponent to $N = 32$, the neutral genealogies.
\end{description}
